\chapter{Introduzione}

\section{Ragionamenti}

Il ragionamento è l'attività fondamentale attraverso cui l'essere umano processa informazioni per giungere a nuove conclusioni.
Nella vita quotidiana ragioniamo continuamente: dalle decisioni più semplici alle inferenze più complesse, utilizziamo schemi di ragionamento che ci permettono di derivare nuova conoscenza a partire da ciò che già sappiamo.
La logica formale studia proprio questi schemi, astraendo dalla particolare natura del contenuto per concentrarsi sulla forma e sulla validità delle inferenze.
Un esempio classico è il \textit{modus ponens}, una regola inferenziale che permette di concludere $B$ a partire dalle premesse $A \rightarrow B$ e $A$: se sappiamo che \qi{se piove, la strada è bagnata} e osserviamo che \qi{piove}, possiamo validamente concludere che \qi{la strada è bagnata}.
In questo caso, la validità del ragionamento non dipende dal contenuto specifico delle proposizioni (pioggia, strada bagnata), ma dalla loro struttura logica: ogni volta che abbiamo un'implicazione e il suo antecedente, possiamo inferire il conseguente.

Non tutti i ragionamenti, tuttavia, seguono la stessa logica, né hanno lo stesso grado di certezza nelle loro conclusioni.
Possiamo distinguere tre forme principali di ragionamento, ciascuna con caratteristiche e applicazioni diverse, che riflettono diversi modi in cui gli esseri umani costruiscono e giustificano la conoscenza.

\begin{definition}{Ragionamento Deduttivo}
  Il \keyword{ragionamento deduttivo} parte da premesse generali per derivare conclusioni necessarie: se le premesse sono vere, la conclusione non può che essere vera.
\end{definition}

È il ragionamento tipico della matematica e della logica formale, dove la validità è garantita dalla sola struttura dell'argomento.
In un certo senso, il ragionamento deduttivo non aggiunge nuova informazione, ma piuttosto esplicita ciò che è già implicito nelle premesse.

\begin{definition}{Ragionamento Abduttivo}
  Il \keyword{ragionamento abduttivo} procede dall'effetto alla causa, cercando la spiegazione più plausibile per un fenomeno osservato.
\end{definition}

È il ragionamento tipico della diagnosi medica o dell'investigazione scientifica, dove si osservano dei dati e si cerca di inferire la causa più probabile, pur sapendo che potrebbero esserci altre spiegazioni possibili.

\begin{definition}{Ragionamento Induttivo}
  Il \keyword{ragionamento induttivo} procede dal particolare al generale: osservando un certo numero di casi, si induce una regola o una legge generale.
\end{definition}
Qui la conclusione è probabile ma non certa, e la forza dell'inferenza dipende dal numero e dalla qualità delle osservazioni.

Ognuno di questi tipi di ragionamento ha un ruolo fondamentale nella nostra comprensione del mondo e nella nostra capacità di agire efficacemente e per ognuno di essi esistono formulazioni diverse del concetto di validità e di conseguenza logica.
In questo corso ci concentreremo principalmente sul ragionamento deduttivo e sui sistemi formali che ne catturano la struttura. Studieremo come rappresentare il linguaggio naturale attraverso linguaggi formali, quali la \keyword{logica proposizionale} e la \keyword{logica dei predicati}, come definire rigorosamente la nozione di conseguenza logica, e come costruire sistemi di deduzione che ci permettano di derivare meccanicamente conclusioni valide.

\section{Validità e Conseguenza Logica}

Abbiamo visto che esistono diverse forme di ragionamento, ma come possiamo stabilire quando un ragionamento è corretto?
Nel contesto della logica deduttiva, la nozione fondamentale è quella di \textit{validità}: un ragionamento valido è tale che, se le premesse sono vere, la conclusione non può essere falsa.
Questa proprietà dipende unicamente dalla struttura logica dell'argomento, non dal contenuto specifico delle proposizioni coinvolte.

\begin{definition}{Ragionamento valido (informalmente)}
  Un ragionamento è detto \keyword{valido} se, in ogni situazione in cui le premesse sono tutte vere, anche la conclusione è vera
\end{definition}

La validità è strettamente connessa a un altro concetto fondamentale: quello di consistenza (o coerenza) di un insieme di enunciati.
Intuitivamente, un insieme di affermazioni è consistente se è possibile che siano tutte vere contemporaneamente, mentre è inconsistente se contiene una contraddizione che rende impossibile la loro verità simultanea.

\begin{definition}{Inconsistenza (informalmente)}
  Un insieme di enunciati è detto \keyword{inconsistente} se non esiste nessuna \qi{situazione} immaginabile in cui gli enunciati sono tutti veri
\end{definition}

È importante notare che l'inconsistenza è una proprietà strutturale, non empirica: dipende dalle relazioni logiche tra gli enunciati, non dal fatto che siano effettivamente veri o falsi nel mondo reale.

\begin{note}{}
  L'inconsistenza non riguarda la verità o falsità degli enunciati, ma la possibilità di rendere tutti gli enunciati veri contemporaneamente.
  Anche un insieme di enunciati tutti falsi può non essere inconsistente, se esiste una situazione in cui tutti gli enunciati sono veri.
\end{note}

Il concetto duale all'inconsistenza è naturalmente quello di consistenza, che caratterizza gli insiemi di enunciati privi di contraddizioni logiche.

\begin{definition}{Consistenza (informalmente)}
  Un insieme di enunciati è detto \keyword{consistente} se esiste almeno una \qi{situazione} immaginabile in cui gli enunciati sono tutti veri
\end{definition}

\begin{note}{}
  Essere consistenti significa non essere inconsistente.
\end{note}

Per comprendere meglio questi concetti, consideriamo un esempio concreto che illustra come un insieme di enunciati apparentemente innocui possa in realtà essere inconsistente.

\begin{example}{}
  Consideriamo l'insieme di enunciati:
  \begin{enumerate}
    \item O Carlo non dorme, o Carlo sogna\label{ex_item:inconsistency_1}
    \item Carlo dorme\label{ex_item:inconsistency_2}
    \item Se Carlo sogna allora Carlo è felice\label{ex_item:inconsistency_3}
    \item Carlo non è felice\label{ex_item:inconsistency_4}
  \end{enumerate}

  Questi enunciati sono chiaramente inconsistenti. Infatti, dato che Carlo dorme (\ref{ex_item:inconsistency_2}), deve necessariamente sognare (\ref{ex_item:inconsistency_1}). 
  Ma se Carlo sogna, allora deve essere felice (\ref{ex_item:inconsistency_3}), il che contraddice l'enunciato (\ref{ex_item:inconsistency_4}) secondo cui Carlo non è felice.
  Abbiamo quindi una contraddizione che rende l'insieme inconsistente.
\end{example}

I concetti di validità e consistenza che abbiamo introdotto ci permettono ora di definire in modo preciso la nozione centrale della logica: quella di \textit{conseguenza logica}.


Un enunciato è conseguenza logica di un insieme di premesse se segue necessariamente da esse, ovvero se non è possibile che le premesse siano tutte vere e la conclusione falsa.

Questa definizione può essere formulata elegantemente in termini di inconsistenza: un enunciato segue logicamente da un insieme di premesse se assumere la sua negazione, insieme alle premesse, porta a una contraddizione.

\begin{definition}{Conseguenza Logica}
  Dato un insieme \(\Gamma\) di enunciati e un enunciato \(\phi\), diremo che \(\phi\) è \keyword{conseguenza logica} di \(\Gamma\) se l'insieme \(\Gamma \cup \{\neg \phi\}\) è inconsistente. % NDR: Excluded middle assumed
\end{definition}

Si noti che, se $\phi$ è conseguenza logica di $\Gamma$, allora, secondo la definizione: se tutti gli enunciati in $\Gamma$ sono veri, allora $\neg \phi$ è necessariamente falso, ovvero $\phi$ è vero; se $\neg \phi$ è vero, allora almeno uno degli enunciati di $\Gamma$ dovrà essere falso. Il primo caso cattura l'idea della dimostrazione diretta, mentre il secondo caso cattura l'idea della dimostrazione per assurdo.
In particolare, vediamo come il ragionamento per assurdo si applica ad un caso concreto.

\begin{example}{}
  Dati gli enunciati \(\Gamma\):
  \begin{enumerate}
    \item Se Carlo è italiano e Giovani non è spagnolo, allora Rachele è tedesca\label{ex_item:logical_consequence_1}
    \item Se Rachele è tedesca allora Lucia è francese oppure Giovanni è spagnolo\label{ex_item:logical_consequence_2}
    \item Se Lucia non è francese allora Carlo è italiano\label{ex_item:logical_consequence_3}
    \item Giovanni non è spagnolo\label{ex_item:logical_consequence_4}
  \end{enumerate}
  L'enunciato \(\phi\):
  \begin{center}
    Lucia è Francese
  \end{center}
  è conseguenza logica di \(\Gamma\)?
  
  Per verificarlo, dobbiamo valutare se l'insieme \(\Gamma \cup \{\neg \phi\}\) è inconsistente.
  In altre parole, procediamo con una dimostrazione per assurdo.
  
  Assumiamo che Lucia non sia francese (\(\neg \phi\)).

  Allora, per l'enunciato (\ref{ex_item:logical_consequence_3}), Carlo è italiano. 

  Dato che Giovanni non è spagnolo (\ref{ex_item:logical_consequence_4}) e Carlo è italiano, dall'enunciato (\ref{ex_item:logical_consequence_1}) concludiamo che Rachele è tedesca.

  Ma se Rachele è tedesca, per l'enunciato (\ref{ex_item:logical_consequence_2}), Lucia è francese oppure Giovanni è spagnolo.
  Poiché Giovanni non è spagnolo (\ref{ex_item:logical_consequence_4}), deve necessariamente essere vero che Lucia è francese.
  Questo contraddice la nostra ipotesi iniziale che Lucia non sia francese.
  Abbiamo quindi ottenuto una contraddizione, il che dimostra che \(\Gamma \cup \{\neg \phi\}\) è inconsistente, e quindi \(\phi\) è conseguenza logica di \(\Gamma\).
\end{example}

Andiamo dunque a definire formalmente la nozione di \keyword{logica}, per andare a formalizzare i concetti visti finora in modo rigoroso e sistematico.

\begin{definition}{Logica}
  Una \keyword{logica} è un sistema formale caratterizzato da:
  \begin{itemize}
    \item Un insieme di simboli e regole sintattiche che definiscono un \keyword{linguaggio} formale, che ne delinea la sintassi e la struttura delle espressioni;
    \item Una \keyword{semantica}, che assegna significati alle espressioni del linguaggio;
    \item Un \keyword{calcolo}, che definisce le regole di inferenza.
  \end{itemize}

  Il \keyword{calcolo} funge in qualche modo da ponte tra la sintassi e la semantica, essendo un insieme di regole di manipolazione simbolica (trasformazioni sintattiche) che preservano certe proprietà semantiche, come la verità o la validità.
\end{definition}

Esempi di logiche sono:
\begin{itemize}
  \item Logica Proposizionale
  \item Logiche Modali (Tra cui la logica temporale)
  \item Logica del Primo Ordine
\end{itemize}

Durante il corso, ci concentreremo sulle logiche dette \textit{classiche}, che si basano su principi come il \textit{terzo escluso} e il \textit{non contraddittorio}, e che sono caratterizzate da una semantica bivalente (ogni proposizione è vera o falsa).
Tra queste in particolare, studieremo la logica proposizionale e la logica del primo ordine, che sono le fondamenta su cui si basano molte altre logiche più complesse.